\section{The comparison through every deformation jet}\label{sec:jets}

The preceding section supplied the actual completed homotopy terms and a
finite abutment filtration. We now identify their coefficient cochains
through every quotient \(k[[x]]/(x^m)\). The first step lifts the marked
split \(p\)-divisible group. The second recovers the formal law from its
finite torsion, evaluates ordinary cooperations, and passes to the inverse
limit with its actual semilinear action.

\subsection{Lifting the marked split deformation}

The next lemma provides an actual coefficient map through every
$x$-jet.  We continue to use the algebraic $p$-divisible group
$\mathcal G$ over $A=k[[a,t]]$, whose finite kernels were formed
before passage to $K$. In the equivariance statement below, $P$ acts
on the source ring $A$ by its height-three deformation action, whereas
$J$ fixes $A$. The actions on the target include the marking and
height-two deformation actions specified in the statement.

\begin{lemma}[Lifting the split deformation]\label{lem:split-jet-lift}
There are compatible maps of $k$-algebras
\begin{equation}\label{eq:split-parameter-map}
  \varphi_m:A\longrightarrow R[x]/(x^m),\qquad m\geq1,
\end{equation}
reducing to $a\mapsto0$, $t\mapsto t$ at $x=0$, with the following
property.  The pulled-back $p$-divisible group is isomorphic, with its
specified identification at $x=0$, to the direct sum of the universal
height-two equal-characteristic deformation modulo $x^m$ and the
constant height-one \'{e}tale group.

The maps and isomorphisms respect $P$, acting on $R$ and fixing $x$,
and $J$, acting both on $R$ by marking changes and on $k[[x]]$ by the
height-two deformation action.  For fixed $m$, the parameter map is
continuous into one complete finite root/marking stage with nilpotent
variable $x$.  Any finite set of coefficients of the formal-group
isomorphism lies in such a stage after a further finite enlargement.
\end{lemma}

\begin{proof}
\emph{The split special fiber.}
The two markings of Lemma~\ref{lem:etale-marking} identify the
connected part of $\mathcal G_R$ with $\Gamma_2$ and its
\'{e}tale quotient with $\Qp/\Zp$.
Their extension splits canonically over $R$.  This also follows
directly from Lemma~\ref{lem:lift-torsors}: the compatible inclusions
$L^{1/p^{2l}}\to R$ give the unique lifts of the marked
\'{e}tale generators over the perfect ring.  Thus
\[
  \mathcal G_R\simeq\Gamma_2\oplus\Qp/\Zp.
\]
Let $\mathcal H_m$ be the universal height-two deformation over
$A_m$, pulled back to $R[x]/x^m$, and put
\[
  \mathcal D_m=\mathcal H_m\oplus\Qp/\Zp.
\]
The displayed splitting specifies the initial identification with
$\mathcal G_R$.

\smallskip\noindent\emph{Square-zero deformation theory.}
The parameter map used here is nonlocal: it sends $t$ to a unit. We use
square-zero deformation theory formulated for arbitrary base ring maps.
This is Grothendieck's lifting theory for Barsotti--Tate groups, as
presented in \cite[Theorem~4.4 and Corollary~4.7]{Illusie85}; the Hodge
filtration formulation for locally liftable groups appears in
\cite[Chapter~V, Theorem~(1.6)]{Messing72}. The precise square-zero
interface used below is Lau's formulation.

We recall precisely the deformation-theoretic input.  For a
$p$-divisible group $G$ over a ring on which $p$ is nilpotent, isomorphism
classes of lifts across a square-zero ideal $I$ form a torsor under
\[
  \Hom(\omega_{G^\vee},I\otimes\Lie G).
\]
Changing a lifting ring map acts through the Kodaira--Spencer
homomorphism \cite[Theorem~5.1 and Equation~(5.1)]{Lau10}.
For the universal deformation \(\mathcal G/A\), where \(A=k[[a,t]]\)
is a complete local Noetherian \(k\)-algebra with residue field \(k\)
and \(k\) is perfect, we use Lau's separate universal-deformation
statement
\cite[paragraph following Equation~(5.2), p.~227]{Lau10}.
As in Equation~\eqref{eq:KS-decompletion}, its closed-point criterion
and Lemma~\ref{lem:finite-p-basis} give
\[
 \kappa_{\mathcal G}:
 \Hom_A(\Lie\mathcal G,\omega_{\mathcal G^\vee})
 \xrightarrow{\ \sim\ }\Omega^1_{A/\mathbb Z}
 \simeq\Omega^1_{A/k}=A\,da\oplus A\,dt.
\]
At a square-zero jet stage, the ideal \(I\) is an \(A\)-module through
\(\varphi_m\). Applying \(\Hom_A(-,I)\) to this isomorphism over \(A\)
gives the isomorphism of torsor-action groups used in Lau's Equation~(5.1).

\smallskip\noindent\emph{Parameter lifts and their uniqueness.}
We must first exhibit ring lifts for the nonlocal specialization.
For nilpotent corrections, prescribe
\[
  a\longmapsto\alpha,\qquad
  t\longmapsto t+\beta,\qquad
  \alpha,\beta\in xR[x]/x^m.
\]
For $f\in A=k[[a,t]]$, define its Hasse derivatives using independent
formal variables $z,w$ by
\[
 f(a+z,t+w)=\sum_{r,s\geq0}
   (\partial_a^{[r]}\partial_t^{[s]}f)(a,t)\,z^r w^s
 \quad\text{in }k[[a,t,z,w]]=A[[z,w]].
\]
These prescriptions extend uniquely to a ring map by the finite
Hasse--Taylor expression
\begin{equation}\label{eq:finite-taylor}
 f\longmapsto
   \sum_{\substack{r,s\geq0\\r+s<m}}
     (\partial_a^{[r]}\partial_t^{[s]}f)(0,t)\,
        \alpha^r\beta^s.
\end{equation}
One may obtain the same formula without topology by expanding $A$
over $A^{p^e}$ in the finitely many monomials $a^r t^s$ with
$0\leq r,s<p^e$, where $p^e\geq m$; the nilpotent corrections then
disappear on $p^e$-th powers.  This proves both well-definedness and
uniqueness of the substitution. In particular, arbitrary lifts of the
coefficients of $\alpha$ and $\beta$ from one jet to the next give a
ring lift.

Suppose the map and identification have been chosen modulo $x^m$.
The kernel of $R[x]/x^{m+1}\to R[x]/x^m$ is square-zero.
The ring lifts just constructed form a
torsor under the corresponding derivations from $A$.  The map from
this torsor to the torsor of deformations is equivariant for the
Kodaira--Spencer isomorphism.  It is therefore a bijection.
Exactly one parameter lift realizes $\mathcal D_{m+1}$ with the
prescribed identification modulo $x^m$.

There is also at most one isomorphism lifting the prescribed one.
Here the needed rigidity has a short square-zero proof.  If a
homomorphism of $p$-divisible groups reduces to zero across a
square-zero ideal in characteristic $p$, its values lie in the
infinitesimal kernel at the identity.  That kernel is an additive
group of derivations into the square-zero ideal and is killed by
$p$.  Indeed, on any test algebra and at each finite torsion level,
a point reducing to the identity has the form $\varepsilon+d$,
where $\varepsilon$ is the identity point and $d$ is such a
derivation; the group law adds these derivations.  The
homomorphism is consequently killed by $p$; since multiplication by
$p$ on its source $p$-divisible group is an epimorphism for the flat
topology, the homomorphism is zero.  Applying this to the difference
of two lifts proves uniqueness.  Induction proves compatibility of
all maps and isomorphisms in \eqref{eq:split-parameter-map}.

\smallskip\noindent\emph{Compatibility with the two stabilizer actions.}
The unique lifts must respect both the action of $P$ on the deformation
and the marking-change action of $J$. We check the special-fiber action
and then use uniqueness at each square-zero step. On the perfect special
fiber there are
no cross homomorphisms between the marked connected and
\'{e}tale summands.  A section of connected Honda torsion over
the reduced ring $R$ is zero, and a map in the opposite direction
vanishes by connectedness.  The transported action is therefore
block diagonal.  By Lemma~\ref{lem:etale-marking}, $P$ preserves the
connected frame and changes the \'{e}tale frame by a constant
norm unit.  The action of $J$ changes the connected frame by a
constant Honda automorphism and the \'{e}tale frame by a
constant unit.  These actions extend to $\mathcal D_m$: use the
universal height-two deformation action on its connected factor and
the same units on the constant factor.  

More explicitly, let $f:\mathcal G_R^0\to\Gamma_2$ be the connected
marking and let $s_f:\Gamma_2\oplus\Qp/\Zp\to\mathcal G_R$ be the
splitting supplied by $f$ and the uniquely lifted normalized
generator $e_f$.  In the convention of
\eqref{eq:normalized-frame-action},
\[
 f'=jfg^{-1},\qquad
 s_{f'}=g s_f
       \bigl(j^{-1}\oplus[\Nrd(j)\Nrd(g)^{-1}]\bigr).
\]
The equality is checked on the connected marking and on the
normalized \'{e}tale generator, which determine both summands.
This is the block action that we lift to $\mathcal D_m$, using
the universal height-two action and the same constant unit on
the \'{e}tale factor.

Functoriality of the two
lifting torsors, followed by uniqueness, makes
\eqref{eq:split-parameter-map} and its group isomorphism equivariant.

\smallskip\noindent\emph{Control at a finite coefficient stage.}
Write $\varphi_m(a)=\alpha$ and $\varphi_m(t)=t+\beta$ for the selected
map. For fixed $m$, the finitely many coefficients of $\alpha$ and $\beta$
belong to a common finite root/marking stage.  The Hasse derivatives
in \eqref{eq:finite-taylor} are elements of $k[[t]]$ and depend
continuously on $f$ in its $(a,t)$-adic topology.  For fixed $m$ their
orders in $t$ tend to infinity with the $(a,t)$-adic order of $f$,
up to a fixed finite shift.  Formula~\eqref{eq:finite-taylor}
therefore gives continuity into that complete stage.  Every
individual coefficient of the formal-group isomorphism is an
element of $R[x]/x^m$.  A finite set of them lies in a common larger
stage.  This last assertion is deliberately about finite sets of
coefficients; no single finite stage for the whole marking is
required.
\end{proof}

\subsection{Comparison of the homotopy cochains}

\begin{theorem}[The comparison through all jets]
\label{thm:jet-comparison}
For every $m\geq1$ there is a natural, compatible, $J$-equivariant
quasi-isomorphism of cochain complexes
\begin{equation}\label{eq:finite-jet-map}
  D_m^\bullet\longrightarrow\Ccts(P,R)[x]/x^m,
\end{equation}
where the cochains on $R$ have the finite-stage convention of
Definition~\ref{def:stage-cochains}.  The action of $P$ on the target
fixes $x$, and $J$ acts through its actions on $R$ and the height-two
coefficients.  Consequently, in each internal degree $2j$, the
diagonal homotopy cochains of the spectral sequence in
Proposition~\ref{prop:descent-test} are naturally quasi-isomorphic to
\begin{equation}\label{eq:all-jet-comparison}
  \left(\lim_m\Ccts(P,R)[x]/x^m\right)
            \otimes_{k[[x]]}\Omega^{\otimes j}.
\end{equation}
Odd internal degrees vanish.  In particular,
\begin{equation}\label{eq:homotopy-second-page}
  \mathsf E_{2,\Delta}^{s,2j}
      =H^s(P,R)[[x]]\otimes_{k[[x]]}\Omega^{\otimes j}.
\end{equation}
Here the coefficient groups are the actual representations of
Theorem~\ref{thm:coefficient-calculation}: they are $k$ in degrees
$0,1$, the one-dimensional line $\ell$ in degrees $3,4$, and zero
otherwise.  The action of $J$ on the first two lines is trivial, and
$\ell|_{\mathcal O_F^\times}=k(\delta)$.
\end{theorem}

\begin{proof}
\emph{Recovering the formal isomorphism.}
Apply Lemma~\ref{lem:split-jet-lift}.  We first show directly that
its isomorphism of algebraic $p$-divisible groups supplies the
formal-law isomorphism required by ordinary cooperations.
Fix $m$, put $\Lambda=R[x]/x^m$, and write $I=(x)$.
Let $\mathcal F_3$ be the formal law over $A$.  The preparation used to form
the finite kernel over $A$ is
\[
 [p^l]_{\mathcal F_3}(T)=U_l(T)W_l(T),\qquad U_l\in A[[T]]^\times,
\]
where $W_l$ is monic.  Set
\[
 \mathcal B_l=\Lambda[T]/\varphi_m(W_l),\qquad
 f_l(T)=\varphi_m([p^l]_{\mathcal F_3}(T))\in\Lambda[[T]].
\]
Coefficientwise application of $\varphi_m$ preserves the displayed
identity.  Moreover $\varphi_m(U_l)$ is still a unit, because its
constant coefficient is the image of an $A$-unit.  Thus
$(f_l)=(\varphi_m(W_l))$ in $\Lambda[[T]]$.

Put $n_l=p^{2l}$.  Modulo $I$, the law has height two and
\[
              \overline f_l=T^{n_l}v_l(T),
                  \qquad v_l(0)\in R^\times.
\]
In $Q_l=\Lambda[[T]]/(f_l)$ this gives $T^{n_l}\in IQ_l$;
hence $T^{mn_l}=0$, since $I^m=0$.  In the other direction,
the invariant-differential formula in characteristic $p$ gives
$[p]'(T)=0$.  The series $[p](T)$ therefore factors through
$T^p$, and its $l$-fold composite belongs to $(T^{p^l})$.
Consequently
\begin{equation}\label{eq:torsion-coordinate-bounds}
        (T^{mp^{2l}})\subseteq(f_l)\subseteq(T^{p^l}).
\end{equation}
These nested ideals are cofinal with the powers of $(T)$.

We also identify the finite algebraic factor explicitly.  The
reduced prepared polynomial has the factorization
\[
 \overline{\varphi_m(W_l)}=T^{n_l}h_l(T),
                  \qquad h_l(0)\in R^\times.
\]
Its factors are relatively prime, so the Chinese remainder theorem
gives
\[
 \mathcal B_l/I\mathcal B_l
       \cong R[T]/T^{n_l}\times R[T]/h_l.
\]
The identity idempotent lifts uniquely across the nilpotent ideal
$I\mathcal B_l$, giving
\[
          \mathcal B_l=\mathcal B_l^0
                         \times\mathcal B_l^{\mathrm{away}}.
\]
On $\mathcal B_l^0$, $T^{n_l}\in I\mathcal B_l^0$ and thus
$T^{mn_l}=0$.  On $\mathcal B_l^{\mathrm{away}}$, $T$ is
invertible modulo $I$ and hence invertible.  The polynomial map
$\mathcal B_l\to Q_l$ kills the away factor, because $T$ is
nilpotent in $Q_l$.  Conversely, evaluation at the nilpotent
element $T\in\mathcal B_l^0$ defines a map
$Q_l\to\mathcal B_l^0$.  These two maps are inverse on coefficients
and $T$; both algebras are generated by polynomials in $T$.
Therefore
\[
       \Lambda[[T]]/(f_l)\cong\mathcal B_l^0.
\]
The $T$-adic completion of $\mathcal B_l$ stabilizes at this factor.
In particular this identification uses no Noetherian hypothesis on
$\Lambda$.

The factor $\mathcal B_l^0$ is the relative identity component of
the finite kernel.  Its unique idempotent construction is natural
under base change and frame isomorphisms.  The isomorphism of
$p$-divisible groups therefore restricts compatibly to these
factors.  By \eqref{eq:torsion-coordinate-bounds}, their inverse
limit recovers $\Lambda[[T]]$. To recover the group law as well, put
$N_l=mp^{2l}$. The same bounds give
\[
 (T,S)^{2N_l-1}
   \subseteq(f_l(T),f_l(S))
   \subseteq(T,S)^{p^l}
       \quad\text{in }\Lambda[[T,S]].
\]
Thus every finite order in the two coordinates is seen at a sufficiently
large common torsion level. The compatible homomorphism identities on
the products of the finite kernels determine the identity on the full
formal laws. We have obtained the required continuous formal-law
isomorphism and its inverse, naturally under all the marking actions.

\smallskip\noindent\emph{Evaluating the cooperation complex.}
This isomorphism evaluates the degree-zero cooperation algebra in
$R[x]/x^m$.  To include the other factor in
\eqref{eq:ordinary-kunneth}, apply $\varphi_m$ pointwise to its
characteristic-$p$ height-three continuous functions.  Formula
\eqref{eq:finite-taylor} makes this a continuous evaluation into a
complete finite root/marking stage.  More explicitly, an element of
the ordinary tensor product is a finite sum of cooperation elements
times continuous coefficient functions.  For this element and this
$m$, choose a stage containing the parameter corrections and the
finitely many formal-isomorphism coefficients occurring in those
cooperation elements.  All its evaluated functions then take values
in that one stage.  This proves the finite-stage requirement for
\eqref{eq:finite-jet-map}; it does not demand one stage for all
cooperation elements at once.  The same argument retains arbitrary
compact profinite parameters, using the same finite Taylor formula
and supremum topology.

\smallskip\noindent\emph{Agreement with the realized stabilizer action.}
To use this coefficient map in the homotopy test, we must identify its
\(J\)-action with the action induced by the actual automorphisms of
\(E_2\). This requires the full marking isomorphism through every jet.
We make that agreement explicit. Let
\(\mathcal F_2\) be the chosen height-two law over
\(B_2=W(k)[[x]]\). For \(j\in J\), write its marked-lift transport as
\[
 \rho_j:B_2\longrightarrow B_2,\qquad
 h_j:\rho_j^*\mathcal F_2\xrightarrow{\sim}\mathcal F_2,\qquad
 h_j\bmod(p,x)=j.
\]
The last equality retains the entire Honda series. These are the
coefficient action and universal transport of
\citet[Section~1, Equation~(1.4), pp.~673--674]{DH95}, in the chosen
coordinate. Uniqueness identifies their reductions with the height-two
action through every jet and gives the group law for the transports.

The geometric marking action \(\mu_j(f)=jf\) induces the left coefficient
action \(\sigma_j^R=(\mu_{j^{-1}})^*\). Let
\(\iota_m:B_2\to\Lambda=R[x]/x^m\) be the coefficient map and
\(\sigma_j\) the diagonal action, so \(\sigma_j\iota_m=\iota_m\rho_j\).
Inverse pullback changes the displayed special-fibre splitting to
\(s_f(j\oplus[\Nrd(j)^{-1}])\). Replacing this connected block by
\(\iota_m(h_j)\), with the same étale unit, restores the same marked
deformation at each lifting step. The two uniqueness statements in
Lemma~\ref{lem:split-jet-lift} therefore give, for the recovered formal
marking \(\Phi_m:\varphi_m^*\mathcal F_3\to\iota_m^*\mathcal F_2\),
\[
 \sigma_j\varphi_m=\varphi_m,\qquad
 \sigma_j^*\Phi_m=\iota_m(h_j^{-1})\circ\Phi_m:
 \varphi_m^*\mathcal F_3\longrightarrow(\iota_m\rho_j)^*\mathcal F_2.
\]
Here the superscript \(\sigma_j^*\) means coefficientwise base change.

\begin{samepage}
For the realized map \(e_j:E_2\to E_2\), let \(\mu\) denote multiplication.
The ordinary same-height cooperation point
\(\pi_*(\mu(1\wedge e_j))\) over \(B_2\) has coefficient maps
\(1,\rho_j\) and right-to-left formal isomorphism \(h_j\)
\citep[Proposition~7.3, Equation~(7.3), and Corollary~7.7]{GH04}.
Thus its coordinate identity is \(X_2=h_j(e_j(X_2))\). Let
\(F_{E_3}/E_{3,0}\) denote the chosen height-three law before reduction.
A degree-zero point of the ordinary cross-height algebra over a commutative
ring \(Q\) is a triple \((\psi_2,\psi_3,\Phi)\):
\(\psi_2:B_2\to Q\) and \(\psi_3:E_{3,0}\to Q\) are its
coefficient maps, and
\(\Phi:\psi_3^*F_{E_3}\xrightarrow{\sim}\psi_2^*\mathcal F_2\)
is its right-to-left formal isomorphism with invertible derivative, as in
Lemma~\ref{lem:ordinary-cooperations}. The actual action of
\(e_j\wedge1\) is consequently
\[
 (\psi_2,\psi_3,\Phi)\longmapsto
       (\psi_2\rho_j,\psi_3,\psi_2(h_j^{-1})\circ\Phi).
\]
\end{samepage}
After diagonal characteristic-\(p\) reduction, this is the same formula
as for \(\Phi_m\), on all isomorphism coefficients; equivalently
\(\psi_2(h_j^{-1})=(\psi_2\rho_j)(h_{j^{-1}})\).
For \(d_j=h_j'(0)\), a compatible degree-two cotangent generator
transforms by \(d_j^{-1}\), and the derivative \(c=u_2/u_3\) of the
marking transforms by the same factor. This identifies the action on
the actual line \(\Omega\): a compatible frame change with linear term
\(a\) replaces \(d_j\) by \(a d_j\rho_j(a)^{-1}\).
The entire series \(j\) is retained, so this includes every principal
unit, even when its closed-fibre derivative is one. Naturality of the
external product in \eqref{eq:ordinary-kunneth} carries this calculation
to every \(C^s\), with \(J\) fixing the height-three function factor.
The bar faces are evaluation, multiplication, and transport by $P$.
Equivariance and uniqueness in Lemma~\ref{lem:split-jet-lift} show
that evaluation commutes with the action face; the remaining faces
and degeneracies commute by their formulas.  The preceding calculation
gives $J$-equivariance.  We have therefore constructed actual maps of
the homotopy cochain complexes in \eqref{eq:finite-jet-map}.
Their construction in nonzero even internal degrees uses only the
height-two periodicity line and hence gives its stated tensor power.

\smallskip\noindent\emph{Quasi-isomorphism at each finite jet.}
At $m=1$, Lemma~\ref{lem:ordinary-cooperations} identifies this map
with
\[
  \Ccts(P,L)\longrightarrow\Ccts(P,R).
\]
It is the natural inclusion and is a quasi-isomorphism by
Theorem~\ref{thm:coefficient-calculation}.  Filter
\eqref{eq:finite-jet-map} by powers of $x$.  Flatness in
Lemma~\ref{lem:ordinary-cooperations} identifies each source
associated-graded layer with $D_1^\bullet$, and each target layer
with $\Ccts(P,R)$.  The induced map is the same residue inclusion,
multiplied by the corresponding formal power of $x$.
The finite filtration thus shows that \eqref{eq:finite-jet-map}
is a quasi-isomorphism for every $m$.

\smallskip\noindent\emph{Passing to the inverse limit.}
Both inverse systems of cochain terms have surjective transition
maps.  Their strict inverse limits consequently compute their
derived inverse limits, and the latter preserve the levelwise
quasi-isomorphisms.  Combining this with
Lemma~\ref{lem:completed-test-terms} proves
\eqref{eq:all-jet-comparison}.  On the target the differential acts
coefficientwise in $x$.  Products of complexes of $k$-vector spaces
are exact, or equivalently the finite-jet cohomology maps are
surjective.  Thus its cohomology is $H^s(P,R)[[x]]$, giving
\eqref{eq:homotopy-second-page} and its asserted coefficient lines.

\smallskip\noindent\emph{The action on completed homotopy.}
We identify the action above with the action on actual completed homotopy.
The original tested spectrum carries $J$, even if the cofibers by $x^m$
were chosen without equivariance. Fix a cosimplicial degree $s$, and put
\[
 D_{\mathrm{ord}}=
   \bigl(\pi_0((E_2\wedge C^s)/p)\bigr)_\Delta,
 \qquad \widehat D=\lim_m D_m^s.
\]
By Lemma~\ref{lem:ordinary-cooperations}, $D_{\mathrm{ord}}$ is flat over
$k[[x]]$ and $D_m^s=D_{\mathrm{ord}}/x^mD_{\mathrm{ord}}$.
Lemma~\ref{lem:completed-test-terms} identifies $\widehat D$ with the
actual degree-zero diagonal completed homotopy and its projections with
reduction maps. For every $m\geq1$,
\[
 \ker(\widehat D\longrightarrow D_m^s)=x^m\widehat D.
\]
Indeed, the component modulo $x^{m+r}$ of an element in this kernel is
$x^m$ times a unique component modulo $x^r$, because multiplication by $x$ is injective on
$D_{\mathrm{ord}}$. These components are compatible as $r$ varies,
proving one inclusion; the reverse inclusion is immediate. The image of $D_{\mathrm{ord}}$ is
dense for the separated topology defined by these projection kernels.
The natural map $(E_2\wedge C^s)/p\to\mathcal T(C^s)$ is $J$-equivariant,
so the actual action agrees with the ordinary action on that image.
For every $g\in J$, semilinearity and $g(x)\in x\,k[[x]]^\times$ give
$g(x^m\widehat D)=x^m\widehat D$. The kernel identity therefore makes
the actual $g$-action continuous for this topology. The action obtained from the
compatible $D_m^s$-actions is also continuous and agrees on the same
dense image, so the two actions coincide. Tensoring with the actual
finite free line $\Omega^{\otimes j}$ gives the same conclusion in every
even internal degree $2j$. All diagonal projections commute with this
action. We have used the new coefficient maps to identify the existing
spectral sequence, without requiring a separate spectral realization
of those maps.
\end{proof}

\begin{remark}\label{rem:why-all-jets}
Theorem~\ref{thm:jet-comparison} identifies the action before
inverting $x$, through every quotient $k[[x]]/(x^m)$.  In particular,
after restriction to $\mathcal O_F^\times$ and the exact scalar
extension to $k((x))$, the upper rows are the
constant character $\delta$ tensored with the actual periodicity
line.  This is the information needed in the following section;
an identification only after setting $x=0$ would not provide it.
\end{remark}
